\documentclass[a4paper,11pt]{article}
\usepackage[utf8]{inputenc}
\usepackage[T1]{fontenc}
\usepackage[french]{babel}
\usepackage{lmodern}
\usepackage{amsmath,amssymb,amsthm}
\usepackage{mathrsfs}
\usepackage{enumitem}
\usepackage{multicol}

\setlist[enumerate]{label=\textbf{\textcolor{Couleur8}{\arabic*.}}, left=1em}

% Si vous voulez aussi modifier les listes enumerate imbriquées :
\setlist[enumerate,2]{label=\textcolor{Couleur8}{\alph*)}}
\setlist[enumerate,3]{label=\textcolor{Couleur8}{\roman*)}}
\setlist[enumerate,4]{label=\textcolor{Couleur8}{\Alph*)}}
\usepackage{graphicx}
\usepackage{xcolor}
\usepackage{tcolorbox}
\usepackage{geometry}
\usepackage{fancyhdr}
\usepackage{titlesec}
\usepackage{cancel}
\usepackage{ifthen}
\usepackage{tikz}
\usetikzlibrary{arrows.meta}
\usepackage{tkz-tab}
\usepackage{eurosym}

% OPTION PROF/ÉLÈVE
% Décommentez UNE des deux lignes suivantes :
\newboolean{prof}\setboolean{prof}{true}   % Version professeur (avec corrections des devoirs)
\newboolean{prof}\setboolean{prof}{true}  % Version élève (sans corrections des devoirs)

\definecolor{Couleur1}{HTML}{4A5D7A} %Th
\definecolor{Couleur2}{HTML}{A5B5D6} %Prop
\definecolor{Couleur3}{HTML}{A8C68A} %Méthode
\definecolor{Couleur6}{HTML}{8A9CC1} %Def
\definecolor{Couleur7}{HTML}{E6B459} %Rq Voc
\definecolor{Couleur8}{HTML}{D36740} %Titres
\definecolor{correction}{HTML}{D36740} % Couleur pour les corrections
\definecolor{coopmathscorr}{HTML}{F15929} % Couleur corrections coopmaths

% Configuration des marges
\geometry{
	top=2cm,
	bottom=2cm,
	inner=1.5cm,
	outer=1.5cm,
}

% Police
\renewcommand{\familydefault}{\sfdefault}

% Compteurs
\newcounter{seancenum}
\newcounter{seancenumD}
\setcounter{seancenum}{1}
\setcounter{seancenumD}{1}

% Configuration des en-têtes et pieds de page
\pagestyle{fancy}
\fancyhf{}
\ifthenelse{\boolean{prof}}{
    \fancyhead[L]{\textcolor{Couleur8}{\textbf{Corrections Automatismes - Classe de seconde - Version Professeur}}}
}{
    \fancyhead[L]{\textcolor{Couleur8}{\textbf{Corrections Devoirs - Séance 13 - Classe de seconde}}}
}
\fancyhead[R]{\textcolor{Couleur8}{\textbf{2025-2026}}}
\fancyfoot[L]{\textcolor{Couleur8}{Mme Bertail et Mme Pinto}}
\fancyfoot[R]{\textcolor{Couleur8}{\thepage}}
\renewcommand{\headrulewidth}{1pt}
\renewcommand{\headrule}{\hbox to\headwidth{\color{Couleur8}\leaders\hrule height \headrulewidth\hfill}}

% Environnements personnalisés
\newtcolorbox{seance}[1]{
    colback=Couleur6!10,
    colframe=Couleur1!65,
    fonttitle=\bfseries\large,
    title=Correction Séance \theseancenum #1,
    left=5mm,
    right=5mm,
    top=3mm,
    bottom=3mm,
    arc=0mm,
    after=\stepcounter{seancenum}
}

\newtcolorbox{seanceD}[1]{
    colback=Couleur6!5,
    colframe=Couleur6!45,
    fonttitle=\bfseries\large,
    title=Correction Devoirs - Séance \theseancenumD #1,
    left=5mm,
    right=5mm,
    top=3mm,
    bottom=3mm,
    arc=0mm,
    after=\stepcounter{seancenumD}
}

\begin{document}

\setcounter{seancenumD}{13}
\newpage
\begin{seanceD}{} %13
\begin{enumerate}
\item  $2{,}1\text{ h} =2 \text{ h} + 0{,}1 \text{ h}$, soit $2\text{ h} + \dfrac{1}{10} \text{ h} = 120 \text{ min} + 6\text{ min} = 126 \text{ min}$.\\
          Ainsi, $2{,}1$ heures correspond à ${\color[HTML]{f15929}\boldsymbol{126}}$ {\color[HTML]{f15929}minutes}.

\item  Dans une $1$ heure, il y a $5\times 12$ minutes.\\
      Ainsi, en $1$ heure, elle parcourt $12$ fois plus de distance, soit $0{,}6\times 12=7{,}2\text{ km}$.\\
      Sa vitesse moyenne est donc ${\color[HTML]{f15929}\boldsymbol{7{,}2}}\text{ km/h}$. 

\item 
 On remplace $x$ par $x-1$ dans l'expression de $m$ :\\
         $\begin{aligned}
        m(x-1)&=-(x-1)+1 \\
        &=-x+1+1\\
        &={\color[HTML]{f15929}\boldsymbol{-x+2}}
        \end{aligned}$


\item  On se ramène à une équation du type $ax=b$ en isolant les  \og{} $x$ \fg{} dans le membre de gauche et les \og{} non $x$ \fg{} dans le membre de droite.\\
    $\begin{aligned}
    (8x-6)-(3x+1)&=0\\
   8x-6-3x-1&=0\\
 5x-7&=0\\
 5x&=7\\
x&= \dfrac{7}{5}
\end{aligned}$\\
    
   La solution de cette équation est  ${\color[HTML]{f15929}\boldsymbol{\dfrac{7}{5}}}$.


\item 
 La fonction $f(x) = 5x+1{,}5$ est une fonction affine de coefficient directeur $a = 5$ et d'ordonnée à l'origine $b = 1.5$.\\
    • La fonction s'annule quand $5x+1{,}5 = 0$, soit $x = \dfrac{-3}{10}$.\\
    • Comme $a = 5 > 0$, la fonction est croissante : elle est négative avant la racine et positive après.\\
    Le tableau de signes de la fonction $f$ est donc le suivant :\\
    \begin{tikzpicture}[baseline, scale=0.75]
    \tkzTabInit[lgt=3,deltacl=0.8,espcl=2.1]{ $x$ / 1.5, $f(x)$ / 1.5}{ $-\infty$, ${-0,3}$, $+\infty$}
	\tkzTabLine{  , -, z, +}
	\end{tikzpicture}

\item 
 En notant $D$ la dépense totale, Alice a dépensé $\dfrac{2}{7}D$ et Louis, $\dfrac{5}{7}D$.\\
Leur participation équilibrée doit être de $\dfrac{1}{2}D$ chacun.\\
Alice doit donc donner à Louis la somme de $\dfrac{1}{2}D - \dfrac{2}{7}D = \dfrac{3}{14}D$.\\
La fraction de la dépense totale que Alice doit donner à Louis est donc ${\color[HTML]{f15929}\boldsymbol{\dfrac{3}{14}}}$.


\end{enumerate}
\end{seanceD}

\end{document}